\section{Per-sampler algorithms}\label{appendix:algorithms}
Algorithm~\ref{alg:unified} in the main text states the shared loop for one assimilation step. Algorithms~\ref{alg:si_sde}--\ref{alg:fm_ode} below specialise it to the three samplers; the differences are confined to the boxed lines.

\begin{algorithm}[ht]
\caption{Sampler 1 -- Stochastic interpolant SDE (native diffusion)}
\label{alg:si_sde}
\begin{algorithmic}[1]
\STATE \textbf{Model:} SI drift $\drift_\theta$; \textbf{anchor} $\bm a_0=\bx_0$; \textbf{diffusion} $\gdiff_\tau=\gamma_\tau$ (native)
\STATE \textbf{Weight} $\gweight_\tau = \vscoef_\tau + \tfrac12\gamma_\tau^2 = \gamma_\tau^2\tau\big(\tfrac{\dot\beta_\tau}{\beta_\tau}\gamma_\tau-\dot\gamma_\tau\big)\big/\gamma_\tau$ \hfill // velocity--score coef.\ for $\sigpath_\tau=\gamma_\tau\sqrt\tau$
\STATE \textbf{Init} $\bXstar\leftarrow\bx_0$ \hfill // point-mass source; reweighting vacuous
\STATE \textbf{Drift} $\drift^{\gdiff}_\tau=\drift_\theta(\bXstar,\bx_0,\tau)$ \hfill // score already built in
\STATE \textbf{Source moments} via Wiener identities (Lemma~\ref{lem:source_moments}, SI case)
\STATE Run Algorithm~\ref{alg:unified} with the Brownian increment \emph{retained}.
\end{algorithmic}
\end{algorithm}

\begin{algorithm}[ht]
\caption{Sampler 2 -- Diffusion-model SDE (lifted diffusion)}
\label{alg:fm_sde}
\begin{algorithmic}[1]
\STATE \textbf{Model:} FM velocity $\fmvel_\theta$; \textbf{anchor} $\bm a_0=\bm 0$, $\bx_0$ as conditioning input; \textbf{diffusion} $\gdiff_\tau>0$ (free, e.g.\ $\gdiff_\tau\propto\sqrt{\alpha_\tau\beta_\tau}$)
\STATE \textbf{Weight} $\gweight_\tau = \vscoef_\tau + \tfrac12\gdiff_\tau^2$, with $\vscoef_\tau=\alpha_\tau(\dot\beta_\tau\alpha_\tau-\dot\alpha_\tau\beta_\tau)/\beta_\tau$
\STATE \textbf{Init} $\bXstar\sim\mathcal{N}(0,\bI)$ \hfill // exact for marginal construction ($\Phi_0^{\obs}$ const)
\STATE \textbf{Score} $\genscore_\tau$ recovered from $\fmvel_\theta$ via Eq.~\eqref{eq:fm_score}; drift $\drift^{\gdiff}_\tau=\fmvel_\theta+\tfrac12\gdiff_\tau^2\genscore_\tau$
\STATE \textbf{Source moments} via Tweedie identities (Lemma~\ref{lem:source_moments}, FM case)
\STATE Run Algorithm~\ref{alg:unified} with the Brownian increment \emph{retained}.
\end{algorithmic}
\end{algorithm}

\begin{algorithm}[ht]
\caption{Sampler 3 -- Guided probability-flow ODE (flow matching, deterministic)}
\label{alg:fm_ode}
\begin{algorithmic}[1]
\STATE \textbf{Model:} FM velocity $\fmvel_\theta$; \textbf{anchor} $\bm a_0=\bm 0$; \textbf{diffusion} $\gdiff_\tau=0$
\STATE \textbf{Weight} $\gweight_\tau = \vscoef_\tau = \alpha_\tau(\dot\beta_\tau\alpha_\tau-\dot\alpha_\tau\beta_\tau)/\beta_\tau$ \hfill // pure velocity--score term
\STATE \textbf{Init} $\bXstar\sim\mathcal{N}(0,\bI)$ \hfill // randomness enters only through the source
\STATE \textbf{Drift} $\drift^{\gdiff}_\tau=\fmvel_\theta(\bXstar,\state^{n-1},\tau)$ \hfill // no score term needed ($\gdiff_\tau=0$)
\STATE \textbf{Source moments} via Tweedie identities (Lemma~\ref{lem:source_moments}, FM case)
\STATE Run Algorithm~\ref{alg:unified} with the Brownian increment \emph{omitted} (deterministic Euler step / higher-order ODE solver).
\end{algorithmic}
\end{algorithm}

\section{Model-specific bookkeeping}\label{appendix:bookkeeping}

\begin{table*}[ht]
    \centering
    \caption{Model-specific quantities for the observation-interpolant method, instantiated for stochastic interpolants and flow matching. The general framework (left column) treats both as affine Gaussian probability paths; the model-specific quantities differ only in the source process $\src_\tau$ and how the velocity and score are obtained. The guided ODE and DM-SDE share the flow matching column; the SI-SDE uses the stochastic interpolant column.}
    \label{tab:si_vs_fm}
    \small
    \begin{tabular}{p{0.235\linewidth}|p{0.345\linewidth}|p{0.345\linewidth}}
        \toprule
        \textbf{Component} & \textbf{Stochastic interpolant} & \textbf{Flow matching} \\
        \midrule
        Interpolant $\bx_\tau$
            & $\alpha_\tau\bx_0 + \beta_\tau\bx_1 + \gamma_\tau\bW_\tau$
            & $\alpha_\tau\latent + \beta_\tau\bx_1,\ \ \latent\sim\mathcal{N}(0,\bI)$ \\
        Source $p_0$
            & $\delta_{\bx_0}$ (previous state) $+$ Brownian
            & $\mathcal{N}(0,\bI)$ (latent), $\bx_0$ as input \\
        Source process $\src_\tau$
            & $\gamma_\tau\bW_\tau$
            & $\alpha_\tau\latent$ \\
        Learned object
            & SDE drift $\drift_\theta=\E[\bR_\tau\mid\bx_\tau]$
            & velocity $\fmvel_\theta=\E[\dot\alpha_\tau\latent+\dot\beta_\tau\bx_1\mid\bx_\tau]$ \\
        Score from model
            & $A_\tau[\beta_\tau\drift_\theta - c_\tau]$ \eqref{eq:SI_score}
            & $\dfrac{\beta_\tau\fmvel_\theta-\dot\beta_\tau\bx}{\alpha_\tau(\dot\beta_\tau\alpha_\tau-\dot\alpha_\tau\beta_\tau)}$ \eqref{eq:fm_score} \\
        Velocity--score coef.\ $\vscoef_\tau$
            & $\dfrac{\sigpath_\tau(\dot\beta_\tau\sigpath_\tau-\dot\sigpath_\tau\beta_\tau)}{\beta_\tau},\ \ \sigpath_\tau=\gamma_\tau\sqrt\tau$
            & $\dfrac{\alpha_\tau(\dot\beta_\tau\alpha_\tau-\dot\alpha_\tau\beta_\tau)}{\beta_\tau}$ \\
        Native generator
            & SDE $\rd\bX=\drift_\theta\rd\tau+\gamma_\tau\rd\bW$
            & ODE $\rd\bX=\fmvel_\theta\rd\tau$ \\
        Available samplers
            & SI-SDE ($\gdiff_\tau=\gamma_\tau$)
            & DM-SDE ($\gdiff_\tau>0$), guided ODE ($\gdiff_\tau=0$) \\
        Observation interpolant $\interpolantobs_\tau$
            & $\alpha_\tau\obsmatrix\bx_0 + \beta_\tau\obs$
            & $\beta_\tau\obs$ \\
        Likelihood mean $\bar\mu_\tau$
            & $\obsmatrix\bx_\tau-\gamma_\tau\obsmatrix\,\E[\bW_\tau\mid\bx_\tau]$
            & $\obsmatrix\bx_\tau+\alpha_\tau^2\obsmatrix\genscore_\tau$ \\
        Source moments
            & Wiener moments (Lemma~\ref{lem:source_moments})
            & Tweedie moments via $\genscore_\tau$ \eqref{eq:fm_source_moments} \\
        Inflated covariance $\bar\Sigma_\tau$
            & $\beta_\tau^2\obscov+\gamma_\tau^2\obsmatrix\Sigma^{\bW}_\tau\obsmatrix^\top$
            & $\beta_\tau^2\obscov+\alpha_\tau^2\obsmatrix(\bI+\alpha_\tau^2\nabla\genscore_\tau)\obsmatrix^\top$ \\
        \bottomrule
    \end{tabular}
\end{table*}

\begin{table*}[ht]
    \centering
    \caption{Approximations and simplifications of the observation-interpolant
    method, grouped by where they enter the pipeline. ``Exact when'' states the
    regime in which the approximation incurs no error. All schedule-derived
    quantities are kept general in $(\alpha_\tau,\beta_\tau,\gamma_\tau)$.}
    \label{tab:approximations}
    \small
    \begin{tabular}{p{0.21\linewidth} p{0.30\linewidth} p{0.20\linewidth} p{0.21\linewidth}}
        \toprule
        \textbf{Approximation} & \textbf{What is assumed / dropped} & \textbf{Where} & \textbf{Exact when} \\
        \midrule
        \multicolumn{4}{l}{\emph{Modelling}}\\
        Linear--Gaussian observations & $\obsoperator$ linear ($\obsmatrix$), noise $\mathcal N(0,\obscov)$ & Lem.~\ref{lem:interpolated_likelihood} & $\obsoperator$ exactly linear \\
        Autoregressive factorisation & Optimal proposal; drop importance weights and future conditioning & Sec.~\ref{sec:preliminaries} & Single assimilation step \\
        \midrule
        \multicolumn{4}{l}{\emph{Likelihood surrogate}}\\
        Observation interpolation & Replace intractable tilt $\Phi_\tau^{\obs}$ by $\interpolantobs_\tau$ surrogate & Sec.~\ref{subsec:obs_interpolation} & $\tau=1$ (and endpoints) \\
        Gaussian likelihood & Interpolant likelihood modelled as Gaussian & Eq.~\eqref{eq:interpolant_likelihood_gaussian} & $p(\bx_1\mid\bx_\tau,\bx_0)$ Gaussian \\
        Fixed-covariance gradients & Differentiate scores with covariances detached & Sec.~\ref{subsec:obs_interpolation} & Locally constant covariance \\
        Inflated covariance (method) & $\bar\Sigma_\tau=\beta_\tau^2\obscov+\obsmatrix\srccov_\tau\obsmatrix^\top$ with full $\srccov_\tau$ & Eq.~\eqref{eq:inflated_covariance} & $p(\bx_1\mid\bx_\tau,\bx_0)$ Gaussian \\
        \midrule
        \multicolumn{4}{l}{\emph{Covariance approximation}}\\
        Jacobian-free covariance & $\srccov_\tau\approx\rho_\tau\bI$; drop network Jacobian & Cor.~\ref{cor:cheap_drift} & $\sigpath_\tau^2\|\nabla^2_{\bx}\log p_\tau\|\ll1$ \\
        Ensemble-shared $\srccov_\tau$ & Full-$\srccov_\tau$ Jacobian evaluated once at the ensemble mean (shared) & Sec.~\ref{subsec:approximations} & Tight ensemble (members coincide) \\
        \midrule
        \multicolumn{4}{l}{\emph{Schedule / numerical}}\\
        Quadratic-$\beta$ SI schedule & $\beta_\tau=\tau^2$; general in $(\alpha,\beta,\gamma)$ & Sec.~\ref{sec:methodology} & Modelling choice (no error) \\
        Endpoint-vanishing $\gdiff_\tau$ & $\gdiff_\tau\!\to\!0$ at endpoints for the DM-SDE lift & Sec.~\ref{subsec:posterior_dynamics} & Keeps FM score finite at $\tau\to1$ \\
        Pseudo-time clamping & Apply correction for $\tau\in[\Delta\tau,1-\Delta\tau]$ & Sec.~\ref{sec:implementation} & $\Delta\tau\to0$ \\
        \bottomrule
    \end{tabular}
\end{table*}

\section{Additional results}\label{appendix:additional_results}

This appendix collects supplementary figures and the extended Navier--Stokes results ($16^2\!\to\!128^2$ super-resolution and $1.5625\%$ sparse scenarios), the ablation study, and the urban-airflow case (data pending).

\begin{figure*}[ht]
    \centering
    \begin{subfigure}{0.24\linewidth}\centering
        \figbox{\linewidth}{2.4cm}{prior $p(x_1|x_0)$}
        \caption{Prior conditional.}\label{fig:an_prior}
    \end{subfigure}\hfill
    \begin{subfigure}{0.24\linewidth}\centering
        \figbox{\linewidth}{2.4cm}{likelihood}
        \caption{Likelihood.}\label{fig:an_like}
    \end{subfigure}\hfill
    \begin{subfigure}{0.24\linewidth}\centering
        \figbox{\linewidth}{2.4cm}{exact posterior}
        \caption{Exact posterior.}\label{fig:an_true}
    \end{subfigure}\hfill
    \begin{subfigure}{0.24\linewidth}\centering
        \figbox{\linewidth}{2.4cm}{sampled posterior}
        \caption{Sampled posterior (one sampler).}\label{fig:an_sampled}
    \end{subfigure}

    \vspace{4pt}
    \begin{subfigure}{0.32\linewidth}\centering
        \figbox{\linewidth}{3.0cm}{KL vs.\ diffusion strength}
        \caption{KL to exact posterior vs.\ diffusion strength $\gdiff_\tau$ (SDE samplers).}\label{fig:an_kl_diff}
    \end{subfigure}\hfill
    \begin{subfigure}{0.32\linewidth}\centering
        \figbox{\linewidth}{3.0cm}{KL vs.\ \# steps}
        \caption{KL to exact posterior vs.\ number of integration steps $M$, per sampler.}\label{fig:an_kl_steps}
    \end{subfigure}\hfill
    \begin{subfigure}{0.32\linewidth}\centering
        \figbox{\linewidth}{3.0cm}{density slices}
        \caption{1-D slices of sampled vs.\ exact posterior densities.}\label{fig:an_slices}
    \end{subfigure}
    \caption{Analytical linear--Gaussian case. Top row: prior conditional, likelihood, exact posterior, and a sampled posterior. Bottom row: convergence of the KL divergence to the exact posterior with respect to the diffusion strength and the number of integration steps, and one-dimensional density slices. \textbf{TODO: insert generated figures.}}
    \label{fig:analytical_panels}
\end{figure*}

\begin{figure*}[ht]
    \centering
    \includegraphics[width=0.92\linewidth]{figures/analytical/analytical_covariance_ablation.pdf}
    \caption{Source-covariance approximation ablation for our three samplers on the analytical linear--Gaussian case: KL to the exact posterior versus the number of sampler steps $M$, for the per-member (exact) Jacobian, the shared ensemble-mean Jacobian, and the isotropic Jacobian-free covariance (Corollary~\ref{cor:cheap_drift}). In this linear--Gaussian case the source covariance is state-independent, so the per-member and shared Jacobians \emph{coincide exactly}; the per-member-vs-shared distinction only matters for a nonlinear prior (the Navier--Stokes and urban cases).}
    \label{fig:analytical_covariance}
\end{figure*}

\begin{figure*}[ht]
    \centering
    \begin{subfigure}{\linewidth}\centering
        \figbox{0.92\linewidth}{1.7cm}{true trajectory snapshots ($t_1,\dots,t_K$)}
        \caption{Ground-truth vorticity over time.}\label{fig:ns_true}
    \end{subfigure}\\[3pt]
    \begin{subfigure}{\linewidth}\centering
        \figbox{0.92\linewidth}{1.7cm}{prior (unconditioned) forecast}
        \caption{Prior forecast (no assimilation).}\label{fig:ns_prior}
    \end{subfigure}\\[3pt]
    \begin{subfigure}{\linewidth}\centering
        \figbox{0.92\linewidth}{1.7cm}{posterior ensemble mean $\pm$ spread}
        \caption{Assimilated posterior (ours).}\label{fig:ns_post}
    \end{subfigure}
    \caption{Navier--Stokes case, one observation scenario. Rows: ground truth, prior forecast, and assimilated posterior. \textbf{TODO: insert generated figures; state observation operator and noise level in caption.}}
    \label{fig:ns_trajectories}
\end{figure*}

\begin{figure*}[ht]
    \centering
    \begin{subfigure}{0.32\linewidth}\centering
        \figbox{\linewidth}{3.0cm}{RMSE vs.\ assimilation step}
        \caption{Ensemble-mean RMSE over a trajectory.}\label{fig:ns_rmse}
    \end{subfigure}\hfill
    \begin{subfigure}{0.32\linewidth}\centering
        \figbox{\linewidth}{3.0cm}{energy spectra}
        \caption{Energy spectra: truth, prior, posterior.}\label{fig:ns_spectrum}
    \end{subfigure}\hfill
    \begin{subfigure}{0.32\linewidth}\centering
        \figbox{\linewidth}{3.0cm}{rank histogram}
        \caption{Rank histogram (calibration).}\label{fig:ns_rank}
    \end{subfigure}
    \caption{Navier--Stokes diagnostics. \textbf{TODO: insert generated figures.}}
    \label{fig:ns_diagnostics}
\end{figure*}

\begin{table*}[ht]
    \centering
    \scriptsize
    \setlength{\tabcolsep}{3.5pt}
    \caption{Navier--Stokes, \emph{Jacobian-free} isotropic covariance (Corollary~\ref{cor:cheap_drift}, \texttt{dps\_jacobian\_free}): ensemble-mean vorticity RMSE, CRPS, and spread--skill ($|1-\mathrm{spread}/\mathrm{skill}|$, $0=$ calibrated) across the four scenarios (super-resolution $16^2\!\to\!128^2$, $32^2\!\to\!128^2$; sparse $5\%$, $1.5625\%$), with per-step cost (lower is better). Companion to the shared-inflated results in Table~\ref{tab:ns_accuracy}; same setup (five held-out trajectories at $E=64$, $M=50$). \tbd{} cells pending the reduced-grid run.}
    \label{tab:ns_accuracy_jacfree}
    \begin{tabular}{l|cccc|cccc|cccc|c}
        \toprule
        & \multicolumn{4}{c|}{\textbf{Vorticity RMSE}}
        & \multicolumn{4}{c|}{\textbf{CRPS}}
        & \multicolumn{4}{c|}{\textbf{Spread--skill}}
        & \textbf{Cost} \\
        & $16^2$ & $32^2$ & $5\%$ & $\tfrac{1}{64}$ & $16^2$ & $32^2$ & $5\%$ & $\tfrac{1}{64}$ & $16^2$ & $32^2$ & $5\%$ & $\tfrac{1}{64}$ & s/step \\
        \midrule
        Ours (SI-SDE) & \tbd & \tbd & \tbd & \tbd & \tbd & \tbd & \tbd & \tbd & \tbd & \tbd & \tbd & \tbd & \tbd \\
        Ours (DM-SDE) & \tbd & \tbd & \tbd & \tbd & \tbd & \tbd & \tbd & \tbd & \tbd & \tbd & \tbd & \tbd & \tbd \\
        Ours (FM-ODE) & \tbd & \tbd & \tbd & \tbd & \tbd & \tbd & \tbd & \tbd & \tbd & \tbd & \tbd & \tbd & \tbd \\
        \bottomrule
    \end{tabular}
\end{table*}

\subsection{Urban airflow}
\label{subsec:results_urban}
The third case is a computational-fluid-dynamics flow over an urban building array: airflow and heat transport through a configuration of bluff bodies (buildings) on a structured domain, with coupled velocity and temperature fields. Relative to the Navier--Stokes case it adds (i) solid obstacles and the associated boundary layers, wakes, and recirculation zones, (ii) a second physical variable (temperature) transported by the flow, and (iii) anisotropic, spatially heterogeneous statistics. It tests the method on a realistic multi-variable applied setting. We use the same two observation scenarios; sensors are placed in physically plausible locations (for the sparse scenario) and the super-resolution scenario observes a coarsened field. Observation noise is additive Gaussian with the per-variable standard deviations stated in the experimental setup. \emph{This case is not yet run: the dataset is still being generated, and the figures and all table cells (\texttt{--}) are placeholders to be filled once the data is available.}

\begin{figure*}[ht]
    \centering
    \begin{subfigure}{0.48\linewidth}\centering
        \figbox{\linewidth}{2.6cm}{velocity: truth / prior / posterior}
        \caption{Velocity field.}\label{fig:urban_vel}
    \end{subfigure}\hfill
    \begin{subfigure}{0.48\linewidth}\centering
        \figbox{\linewidth}{2.6cm}{temperature: truth / prior / posterior}
        \caption{Temperature field.}\label{fig:urban_temp}
    \end{subfigure}
    \caption{Urban airflow case: geometry and reconstructed fields (truth, prior forecast, assimilated posterior) for velocity and temperature. \textbf{TODO: insert generated figures; mark building footprints and sensor locations.}}
    \label{fig:urban_fields}
\end{figure*}

\begin{table*}[ht]
    \centering
    \scriptsize
    \setlength{\tabcolsep}{4pt}
    \caption{Urban airflow, \emph{Jacobian-free} isotropic covariance (Corollary~\ref{cor:cheap_drift}, \texttt{dps\_jacobian\_free}): per-variable ensemble-mean RMSE (velocity $u,v,w$; temperature $\theta_\ell$, fluid cells only), CRPS, and spread--skill ($|1-\mathrm{spread}/\mathrm{skill}|$, $0=$ calibrated) at the two sparse densities, with per-step cost (lower is better). Companion to the shared-inflated results in Table~\ref{tab:urban_accuracy}; same setup (five held-out trajectories at $E=64$, $M=50$). \tbd{} cells pending the reduced-grid run.}
    \label{tab:urban_accuracy_jacfree}
    \begin{tabular}{l|cc|cc|cc|cc|c}
        \toprule
        & \multicolumn{2}{c|}{\textbf{Velocity RMSE}}
        & \multicolumn{2}{c|}{\textbf{Temp.\ RMSE}}
        & \multicolumn{2}{c|}{\textbf{CRPS}}
        & \multicolumn{2}{c|}{\textbf{Spread--skill}}
        & \textbf{Cost} \\
        & $5\%$ & $\tfrac{1}{64}$ & $5\%$ & $\tfrac{1}{64}$ & $5\%$ & $\tfrac{1}{64}$ & $5\%$ & $\tfrac{1}{64}$ & s/step \\
        \midrule
        Ours (SI-SDE) & \tbd & \tbd & \tbd & \tbd & \tbd & \tbd & \tbd & \tbd & \tbd \\
        Ours (DM-SDE) & \tbd & \tbd & \tbd & \tbd & \tbd & \tbd & \tbd & \tbd & \tbd \\
        Ours (FM-ODE) & \tbd & \tbd & \tbd & \tbd & \tbd & \tbd & \tbd & \tbd & \tbd \\
        \bottomrule
    \end{tabular}
\end{table*}

\begin{table*}[ht]
    \centering
    \small
    \caption{Urban airflow: point accuracy per variable, per observation scenario. Lower is better. \emph{Data pending}; all cells (\texttt{--}) are placeholders. There is no ground-truth posterior for the urban case (only a ground-truth state), so KL-at-points is not reported; calibration is assessed by spread--skill and CRPS in Table~\ref{tab:urban_calibration_cost}.}
    \label{tab:urban_accuracy}
    \begin{tabular}{l|cc|cc}
        \toprule
        & \multicolumn{2}{c|}{\textbf{Velocity RMSE}}
        & \multicolumn{2}{c}{\textbf{Temperature RMSE}} \\
        & $5\%$ & $1.5625\%$ & $5\%$ & $1.5625\%$ \\
        \midrule
        Ours (SI-SDE) & -- & -- & -- & -- \\
        Ours (DM-SDE) & -- & -- & -- & -- \\
        Ours (FM-ODE) & -- & -- & -- & -- \\
        \midrule
        FlowDAS & -- & -- & -- & -- \\
        Guided FM (FIG) & -- & -- & -- & -- \\
        Guided FM (OT-ODE) & -- & -- & -- & -- \\
        D-Flow SGLD & -- & -- & -- & -- \\
        SDA & -- & -- & -- & -- \\
        SURGE & -- & -- & -- & -- \\
        \bottomrule
    \end{tabular}
\end{table*}

\begin{table*}[ht]
    \centering
    \small
    \caption{Urban airflow: probabilistic calibration and cost (aggregated over variables unless noted). \emph{Data pending}; all cells (\texttt{--}) are placeholders. No classical (EnKF/PF) baseline: the urban CFD has no true forward solver to propagate the ensemble, so the comparison is generative-only.}
    \label{tab:urban_calibration_cost}
    \begin{tabular}{l|cc|cc|cc}
        \toprule
        & \multicolumn{2}{c|}{\textbf{CRPS}}
        & \multicolumn{2}{c|}{\textbf{Spread--skill}}
        & \multicolumn{2}{c}{\textbf{Cost}} \\
        & $5\%$ & $1.5625\%$ & $5\%$ & $1.5625\%$ & NFE & s/step \\
        \midrule
        Ours (SI-SDE) & -- & -- & -- & -- & -- & -- \\
        Ours (DM-SDE) & -- & -- & -- & -- & -- & -- \\
        Ours (FM-ODE) & -- & -- & -- & -- & -- & -- \\
        \midrule
        FlowDAS & -- & -- & -- & -- & -- & -- \\
        Guided FM (FIG) & -- & -- & -- & -- & -- & -- \\
        Guided FM (OT-ODE) & -- & -- & -- & -- & -- & -- \\
        D-Flow SGLD & -- & -- & -- & -- & -- & -- \\
        SDA & -- & -- & -- & -- & -- & -- \\
        SURGE & -- & -- & -- & -- & -- & -- \\
        \bottomrule
    \end{tabular}
\end{table*}

\subsection{Ablations}
\label{subsec:results_ablations}
On the Navier--Stokes case we isolate the effect of the components introduced in Section~\ref{sec:methodology}. We vary: (i) the covariance along the cost/accuracy axis of Section~\ref{subsec:approximations} -- the Jacobian-free isotropic covariance of Corollary~\ref{cor:cheap_drift} (\texttt{dps\_jacobian\_free}) versus the ensemble-shared inflated source covariance (\texttt{inflated-shared}); (ii) the diffusion strength $\gdiff_\tau$ for the DM-SDE, to map the accuracy/calibration trade-off between the deterministic ODE ($\gdiff_\tau=0$) and strong-diffusion limits; (iii) the number of integration steps $M$; and (iv) the ensemble size $E$. Table~\ref{tab:ablation} collects the results.

\begin{table}[ht]
    \centering
    \small
    \caption{Ablations on the Navier--Stokes case (one observation scenario). Covariance axis (Jacobian-free / inflated-shared) and $\gdiff_\tau$/$M$/$E$ sweeps. \emph{Run in progress}; cells (\texttt{--}) are placeholders, see the expected narrative in the text.}
    \label{tab:ablation}
    \begin{tabular}{l|ccc}
        \toprule
        Configuration & RMSE & CRPS & Spread--skill \\
        \midrule
        Jacobian-free covariance               & -- & -- & -- \\
        Inflated-shared covariance             & -- & -- & -- \\
        \midrule
        $\gdiff_\tau$ sweep (low/med/high) & -- & -- & -- \\
        Steps $M$ (e.g.\ $10/50/100$)      & -- & -- & -- \\
        Ensemble $E$ (e.g.\ $16/64/256$)   & -- & -- & -- \\
        \bottomrule
    \end{tabular}
\end{table}

\paragraph{Analytical: covariance approximation vs.\ the exact posterior.}
On the linear--Gaussian case the exact posterior is available, so we isolate the
likelihood-covariance approximation directly. Table~\ref{tab:analytical_ablation} reports the KL to
the exact posterior for the three choices -- the individual (per-member) Jacobian
(\texttt{inflated}), the ensemble-shared Jacobian (\texttt{inflated-shared}), and the isotropic
Jacobian-free corollary (\texttt{dps\_jacobian\_free}) -- for each of our three samplers. Because the
source covariance is \emph{state-independent} for a linear prior, sharing the Jacobian across the
ensemble introduces no error, so the shared/individual gap only appears for a nonlinear prior (the
Navier--Stokes case). This is the analytical counterpart of the Navier--Stokes covariance ablation
above, with the exact posterior as reference.

\begin{table}[ht]
    \centering
    \small
    \caption{Analytical covariance-approximation ablation: KL divergence to the exact posterior (lower is better), averaged over seeds, per sampler, for the three likelihood-covariance modes. \tbd{} cells pending the reduced-grid run.}
    \label{tab:analytical_ablation}
    \begin{tabular}{l|ccc}
        \toprule
        Covariance mode & SI-SDE & DM-SDE & FM-ODE \\
        \midrule
        Individual Jacobian (\texttt{inflated})        & \tbd & \tbd & \tbd \\
        Shared Jacobian (\texttt{inflated-shared})     & \tbd & \tbd & \tbd \\
        Jacobian-free isotropic (\texttt{dps\_jac.\,free}) & \tbd & \tbd & \tbd \\
        \bottomrule
    \end{tabular}
\end{table}
